\subsection{K8 gains, a capacity counterexample and exact reuse}
The sealed R1 K8 trajectory completes two commits at exact CPD13.228145689936355 ns,2069040 wire,74 SB and5071 CB. It spends160 wire units, saves twelve SB and one CB relative to the same fast R1 root. Its final state quality-dominates the K4 final state by five SB and eight wire units, with CPD/CB equal. The first K8 commit's entire actual route file and fixed endpoint setup/hold vectors exactly equal the two-step K4 final state. Reaching that common physical state uses two versus four candidate FULL/STA gates and five versus nine native children including source/repair observations. Both actors now have final context seals. These exact-state counts substantiate gate amortization; concurrent imported-state timing is not a formal cold complete-algorithm speedup.

The second K8 OFF candidate reaches74 SB and is accepted. ON fails after five of eight restorations at net31767/pin1, restores original occupancy and supplies no candidate route. The controller commits the already validated OFF fallback, following its registered rule. The complete workflow is1057.993 s, including the failed request. Successful-candidate native cores total12.277 s; the failed core is6.423 s, giving18.699 s for all four native requests. The earlier successful-only component summary is preserved with its scope; omitting the failed core from a total would undercount. There are eight native children and three candidate FULL/STA validations, since the incomplete transaction receives no fabricated candidate timing. The complete K8/K4 workflow ratio therefore also reflects different completed validation workloads and must not be advertised as formal acceleration.

\paragraph{An independently verified capacity cause.}
The failed target has a directed path when capacity is ignored, both with and without the retained-crossing domain. Using the native failure's actual transient OCC vector after five successful restorations, the domain has no capacity path, while unrestricted search has an actual-edge witness. The occupancy before any restoration is reconstructed exactly as failure OCC minus the current local NODE multisets plus initial KEEP multisets. Retained-tree inclusion, nonnegative occupancy and add-only restoration are verified. Even in that minimum occupancy, the failed target has no domain-capacity path. Every retained tree node is an optimistic source; native re-expansion and mandatory nonconfigurable closure constraints are relaxed. Absence in this larger search space is a necessary infeasibility certificate for this fixed pruned scope/domain, not a global FPGA routing bound. Because later restorations only add occupied resources, changing their order or increasing weighted-search pressure cannot make this same target/domain feasible. Scope expansion or selective domain relaxation remains possible. Independent Boolean-Floyd oracle checks cover64 topology cases and64 capacity-mask cases; real positive witnesses are checked against actual RRG edges. These diagnostics execute no native rerouting or route flow and do not replace final native legality/STA.

This counterexample changes the next intervention: avoid spending more on the same fixed domain and study per-net guide release or owner-scope expansion. Retain domains that remain useful, let a capacity-incompatible net use the native fallback, then validate the whole combination. It does not establish that releasing only one domain makes the remaining seven jointly feasible or improves quality beyond OFF. The actual OFF witness uses non-retained crossing sources in three selected nets, and the first five ON/OFF partial trees already differ in two nets, so a one-net global-equivalence claim would be unsupported.

\paragraph{A different A0 resource corner.}
The independent resource-only K4 path from pristine strong native A0 finishes at13.607436066820355 ns,2059667 wire,86 SB and5062 CB, spending100 wire units and saving eight SB with unchanged CPD. The timing-then-resource A0 path instead finishes at13.228145689936355 ns,2059623 wire,88 SB and5062 CB. The slower resource-only state has two fewer SB; it remains useful under a tighter SB requirement. The low-SB R1 state spends9417 more wire units and nine more CB than the faster A0 combination while saving fourteen SB. The joint archive now has233 legal observations with separate failures preserved; these are design alternatives under different absolute resource envelopes, not a method-only frontier or per-circuit history-best performance score. Strong A0/R1, U and Tf remain in the population.

\paragraph{Exact cross-batch reuse, with fresh acceptance.}
The serial prototype previously reused candidates only inside one batch. Its registered same-root paired replay contains an OFF request exactly matching a completed native-only transaction. That replay is preserved; a subsequent independent executor adds conservative cross-batch reuse for free native search, requiring the same source paths, bytes, static inputs, objective, request, native binary and pressure. Domain ON and forced paths cannot alias an OFF result. A real completed candidate is reused with zero new native, physical-audit or STA subprocesses; the original envelope accepts it while a zero-wire-headroom policy rejects the identical candidate. Those two local reuse workflows take3.599 and3.302 s and carry no borrowed creation/core time. Domain, forced-path, pressure and request changes miss the cache. A zero-subprocess bookkeeping failure is retained and fixed in the independent validated version; no native experiment was repeated to repair it. These are reuse/interface checks, not a cold speedup.

A new same-root native-only/paired A0 K8 batch is registered behind the existing K4 owner, using the validated executor and completed exact free-native references. It tests whether the R1 gate-amortization effect transfers to the lower-wire A0 start. It retains one native process at a time alongside the two-process main lane, and changes no existing versions/inputs/acceptance rules. Broader resource transfer and native-first controls continue. A separate future in-process candidate/FULL/STA design is motivated by the remaining construction/audit cost; it is not already implemented or timed.
